\par\textbf{Procedimiento completo para comparar tu solución.} El choque normal conserva masa, momento y entalpía total, pero genera entropía. Se calculan razones de presión y densidad, se usa gas ideal para la razón térmica y después se evalúa la diferencia de entropía.

\begin{enumerate}

\item Razón de presiones del choque normal.

\[\fitmath{\mathrm{pr}=1+\frac{2\cdot 1.4}{1.4+1}\cdot \left({2}^{2}-1\right)\approx 4.5\,}\]

\item Razón de densidades por conservación.

\[\fitmath{\mathrm{rr}=\frac{\left(1.4+1\right)\cdot {2}^{2}}{\left(1.4-1\right)\cdot {2}^{2}+2}\approx 2.666667\,}\]

\item La ecuación de estado relaciona ambas razones.

\[\fitmath{\mathrm{Tr}=\frac{\mathrm{pr}}{\mathrm{rr}}}\]

\[\fitmath{\mathrm{Tr}=\frac{(4.5)}{(2.666666667)}\approx 1.6875\,}\]

\item Presión posterior.

\[\fitmath{p_{2}=24000\cdot \mathrm{pr}}\]

\[\fitmath{p_{2}=24000\cdot (4.5)\approx 108000\,\mathrm{Pa}}\]

\item Temperatura posterior.

\[\fitmath{T_{2}=220\cdot \mathrm{Tr}}\]

\[\fitmath{T_{2}=220\cdot (1.6875)\approx 371.25\,\mathrm{K}}\]

\item Mach posterior subsónico.

\[\fitmath{M_{2}=\sqrt{\frac{1+0.2\cdot {2}^{2}}{1.4\cdot {2}^{2}-0.2}}\approx 0.5773503\,}\]

\item Entropía positiva; no usar relaciones isentrópicas a través del choque.

\[\fitmath{\mathrm{ds}=\frac{1.4\cdot 287.05287}{0.4}\cdot \ln\left(\mathrm{Tr}\right)-287.05287\cdot \ln\left(\mathrm{pr}\right)}\]

\[\fitmath{\mathrm{ds}=\frac{1.4\cdot 287.05287}{0.4}\cdot \ln\left((1.6875)\right)-287.05287\cdot \ln\left((4.5)\right)\approx 93.94985\,\mathrm{J/(kg\,K)}}\]

\end{enumerate}
